% ACM v2.20 class and bibliography style are included. See README.md to build.
\documentclass[sigconf,nonacm]{acmart}

% acmart already loads amsmath, booktabs, graphicx, hyperref, and natbib.
% Do not load geometry, titlesec, caption, or a second bibliography system.
\usepackage{mathtools}       % Additional equation tools
\usepackage{bm}              % Bold mathematical symbols
\usepackage{siunitx}         % Aligned numerical columns and units
\usepackage{xcolor}          % Draft notes

% Course milestone: retain ACM typography without publication metadata.
\settopmatter{printacmref=false,printfolios=true}
\setcopyright{none}

% Set to \draftfalse before submission; remove all remaining TODOs first.
\newif\ifdraft
\drafttrue
\ifdraft\emergencystretch=2em\fi
\newcommand{\todo}[1]{\ifdraft\textcolor{red}{[TODO: #1]}\fi}
\newcommand{\method}[1]{\textsc{#1}}
\newcommand{\EM}{\operatorname{EM}}

\title{Reasoning Begins During Reading}
\subtitle{Cross-Worker Evidence Access in Multi-Hop Question Answering}

% Names follow the project README. Confirm spelling and desired order.
% Net IDs are deliberately left blank rather than guessed.
\author{Josh Huang}
\authornote{Net ID: \texttt{JOSH-NETID}.}
\affiliation{\institution{Cornell Tech}\city{New York}\state{NY}\country{USA}}

\author{Elliott Kalt}
\authornote{Net ID: \texttt{esk226}.}
\affiliation{\institution{Cornell Tech}\city{New York}\state{NY}\country{USA}}

\author{Ajit Mallavarapu}
\authornote{Net ID: \texttt{AJIT-NETID}.}
\affiliation{\institution{Cornell Tech}\city{New York}\state{NY}\country{USA}}

\renewcommand{\shortauthors}{Huang, Kalt, and Mallavarapu}

\begin{document}

% ACM places the abstract before \maketitle in the source.
% The title, authors, and abstract appear on the first page, not a cover sheet.
\begin{abstract}
Our working hypothesis is that reasoning begins during reading: models
form relationships between evidence while encoding it, and independently
encoding related documents can make these relationships harder to recover
during answer generation. We investigate this hypothesis using controlled
attention interventions in a two-worker question-answering system.
\todo{Add 2--3 sentences describing the datasets, strongest verified
result, and principal limitation. Keep the abstract around 100--150 words.}
\end{abstract}

\keywords{multi-hop question answering, multi-agent systems, KV cache,
attention masking}
\maketitle

\begin{center}
\small CS 5787 Deep Learning --- Project Milestone --- November 16, 2026
\end{center}

\section{Introduction}
\label{sec:intro}
% Target: 2--3 paragraphs, about 0.4 page.
\todo{Explain the setting: an orchestrator splits evidence across parallel
workers and answers using their outputs. Why is this useful? Why might
independent encoding make composition difficult even without token loss?}

\paragraph{Working hypothesis.}
Reading conditions shape later reasoning. We hypothesize that joint
encoding builds relationships between evidence that are harder to
reconstruct during answer generation.

\paragraph{Research questions.}
\todo{State two questions: (1) Does independent encoding impair composition
when all source tokens are retained? (2) Which cross-worker evidence
access recovers performance, and how does answer-time reasoning interact
with that access?}

\paragraph{Milestone contributions.}
\todo{List only completed contributions: controlled encoding comparison,
paragraph-level attention interventions, and preliminary results.
Distinguish observations from proposed mechanisms.}

\section{Background and Related Work}
\label{sec:related}
% Target: one compact paragraph, about 0.2 page.
\todo{Position the project relative to LatentMAS, Parallel Context Windows,
and Masking in Multi-hop QA. Cite the papers from references.bib after
checking their full text. Explain that your focus is a controlled diagnosis
of evidence access, rather than a demonstrated efficient repair algorithm.}
% Example citation syntax: \cite{pcw2023,masking2025,latentmas2025}.

\section{Problem Formulation and Evaluation}
\label{sec:problem}
% Target: about 0.4 page.
\subsection{Task and Data}
Let an instance be $x=(q,D,y^*)$, where $q$ is a question,
$D=\{d_1,\ldots,d_m\}$ is a set of evidence paragraphs, and $y^*$
is a reference answer. A placement $\pi$ partitions $D$ between two
workers, producing blocks $D_A$ and $D_B$. The orchestrator predicts
$\hat y$ using their encoded evidence.

\todo{Describe synthetic two-fact chains and two-hop MuSiQue, sample
sizes, distractors, and seed. Define AB/BA and co-located placements (AA and BB).
Explain the composition filter and report which
tables use it. Avoid treating mirrored layouts as independent questions.}

\subsection{Metrics and Statistical Analysis}
Our primary metric is normalized exact match:
\begin{equation}
\EM(\hat y,y^*) =
\mathbf{1}\!\left[\operatorname{norm}(\hat y)
=\operatorname{norm}(y^*)\right].
\end{equation}
\todo{Specify normalization, aliases if used, and token F1 if reported.
Average layouts within each source before paired bootstrap confidence
intervals. Separate official scores from exploratory manual same-entity
judgments. State the bootstrap count for each reported experiment.}

\section{Method}
\label{sec:method}
% Target: about 0.5 page.
\subsection{Joint and Independent Encoding}
\todo{Define B1 (joint text encoding), B1-mask (blocked cross-worker
attention at true positions), and K0 (independently encoded, stitched KV
caches). Explain the shared question prefix and position handling.}

\subsection{Controlled Evidence Access}
For a token $i$ in block $B$ and token $j$ in block $A$, define
\begin{equation}
M^{(V)}_{ij} =
\begin{cases}
0, & j\in V,\\
-\infty, & j\notin V,
\end{cases}
\end{equation}
where $M$ is added to attention logits and $V$ is the set of
accessible tokens in $A$. Other attention edges follow the common
causal layout. We vary $V$ between no tokens, the gold supporting
paragraph, a distractor paragraph, all non-gold paragraphs, and all
of $A$. The answer tokens can attend to both blocks in every condition.

\todo{State that the masked runs re-encode B, use oracle paragraph labels,
and do not establish compute savings. Discuss indirect gold-information
relay through A's later tokens when interpreting non-gold controls.}

\subsection{Answer-Time Reasoning}
\todo{Describe thinking-off and forced-thinking conditions, the 1,024-token
thinking budget, answer caps, and recorded cap rates. Thinking-off does
not necessarily mean an absence of reasoning in generated answer text.}

\section{Preliminary Results}
\label{sec:results}
% Target: about 0.6 page. Prefer one primary table and a compact discussion.
\begin{table}[t]
\centering
\caption{Qwen3-8B preliminary masking results, thinking off.
EM (\%), AB/BA pooled, composition-required sources, seed 0.
Values use the existing automatic scorer; scoring limitations remain.}
\label{tab:masking}
\begin{tabular}{lrrrr}
\toprule
Dataset & Sources & Blocked & Gold & Full \\
\midrule
Two-fact toy & 100 & 64.5 & 100.0 & 100.0 \\
MuSiQue & 79 & 14.5 & 38.5 & 29.0 \\
\bottomrule
\end{tabular}
\end{table}

\paragraph{Encoding-time access.}
\todo{Interpret Table~\ref{tab:masking}. Report paired differences and
source-level confidence intervals from the larger masking run. State that
supporting-paragraph access substantially improves accuracy; avoid claiming
that an exact bridge mechanism has been isolated.}

\paragraph{Interaction with answer-time reasoning.}
\todo{Summarize the separate 25-source-per-dataset experiment. Do not merge
its sample with the larger masking table. Discuss cap rates and the latest
scoring audit; do not claim that reasoning adds nothing under gold access.}

\paragraph{Qualitative observations.}
\todo{Optionally give one short example of evidence confusion or report
failure. Label exploratory trace interpretations as hypotheses.}

\section{Limitations and Remaining Work}
\label{sec:next}
% Target: about 0.3 page.
\todo{Cover oracle paragraph selection, one model family, limited seeds,
small/reused samples, shared evidence across sources, scorer sensitivity,
finite generation budgets, and unresolved numerical endpoint differences.}

\todo{Prioritize: verify scoring with a second reader; evaluate a fresh,
checked MuSiQue sample; run clean pre-gold non-gold controls; test a selector
without gold labels. Describe future experiments as planned work.}

\paragraph{Team responsibilities.}
\todo{Assign a concrete completed or planned component to each author.}

\nocite{latentmas2025,pcw2023,masking2025}
\bibliographystyle{ACM-Reference-Format}
\bibliography{references}
\end{document}
